\section{Related Work}

\paragraph{Task-oriented LLM Agents}
LLM agents have been increasingly studied in task-oriented environments that require document understanding~\cite{mmlongbench_doc}, tool use~\cite{Toolformer, ToolLLM}, web interaction~\cite{WebArena, Mind2Web}, or software engineering~\cite{SWE-agent, AutoCodeRover}. A growing body of benchmarks and systems~\cite{AgentBench, OSWorld} evaluates whether agents can follow user instructions, navigate complex environments, and complete prescribed tasks. These settings, however, typically presume that the task has already been specified through a user request, issue description, failing test, or otherwise localized goal, reducing the role of the agent to executing against a stated objective. In contrast, we target the inverse setting in which no such request exists and the relevant problems, often multiple and coexisting, first need to be discovered from a broader context before any of them can be acted on.

\paragraph{Proactive Agents}
Proactive agents aim to move beyond the reactive interaction model by anticipating user needs and initiating assistance before an explicit request is issued. One line of work focuses on uncovering user intent that goes beyond what is literally stated, either by asking clarification questions to resolve ambiguous requests~\cite{DBLP:conf/sigir/AliannejadiZCC19, Kuhn2022CLAMSC, DBLP:conf/emnlp/Zhang0RNC24, DBLP:journals/corr/abs-2511-02208, DiscoverLLM} or by navigating knowledge gaps that have not been articulated~\cite{proper}; these approaches, however, still presume a user-issued query as the anchor of interaction. A more recent line broadens the scope, studying when an agent should intervene~\cite{DBLP:conf/chi/LiuFSWIC25, DBLP:conf/emnlp/ZhangDLMKDCM25}, how user activity or signals can be used to anticipate assistance opportunities~\cite{ProactiveAgent, ContextAgent, FingerTip20K}, and how proactive suggestions should be generated and surfaced~\cite{probe}. Yet across this literature, proactivity remains anchored to a single localized need at a time, leaving open how an agent should jointly surface, ground, and resolve the many coexisting problems that real workflows typically contain.

\paragraph{Templates for LLM Reasoning}
Improving LLM reasoning has traditionally relied on internal model capability, whether by eliciting intermediate steps~\cite{cot, zeroshotcot} or by having a model critique and revise its own outputs~\cite{self-refine, reflexion}. A more recent work observes that useful reasoning patterns recur across problems and externalizes them as reusable templates that can be retrieved and applied: Buffer-of-Thoughts~\cite{BoT} caches prior reasoning traces for retrieval on new problems, and follow-up work extends this idea to hierarchical template paths~\cite{ReasonFlux, SuperCorrect}, to schema-based abstractions for in-context learning~\cite{SA-ICL}, to self-evolving agent memory of reasoning strategies~\cite{ReasoningBank}, to graph-based reuse of thought fragments~\cite{RoT}, and to multi-hop reasoning over long-context documents~\cite{ToTAL}. These template-based approaches, however, share a common assumption that the problem statement is already given and templates serve as schemas for how to solve it. We instead repurpose templates as discovery schemas that specify what contextual signals to attend to and how to connect them in order to infer a problem that has not been stated, and apply them iteratively so that each round extends coverage over coexisting problems rather than refines a single solution.
